% Mathematical extension for the thesis chapter on hemodynamically constrained
% attention. Self-contained: \input this into a document with amsmath, amsthm, amssymb.
%
% Every numerical value quoted in the remarks is produced by
%   python experiments/aggregate.py --readme README.md
% and appears in the generated block of README.md. Nothing here is hand-entered.

\section{Conserved supply as a gain field on attention heads}

\subsection{Notation}

We follow the matrix conventions of Kozachkov, Kastanenka and Krotov
\cite{kozachkov2023}. Let $X \in \mathbb{R}^{d \times N}$ collect $N$ token
embeddings of dimension $d$, and let a cross-attention layer with $H$ heads and
head width $d_k$ produce, for head $h \in [H]$,
\begin{equation}
  Q_h = W_Q^{(h)} X, \qquad K_h = W_K^{(h)} Y, \qquad V_h = W_V^{(h)} Y,
  \qquad
  A_h = \operatorname{softmax}\!\left( \frac{Q_h^{\top} K_h}{\sqrt{d_k}} \right),
\end{equation}
with per-head output $O_h = V_h A_h^{\top} \in \mathbb{R}^{d_k \times N}$. Write
$W_O^{(h)} \in \mathbb{R}^{d_{\mathrm{out}} \times d_k}$ for the $h$-th column block
of the output projection, so that the unconstrained layer computes
\begin{equation}
  \label{eq:dense}
  \mathcal{T}(X) \;=\; \sum_{h=1}^{H} W_O^{(h)} O_h .
\end{equation}

\subsection{The supply operator}

\begin{definition}[Conserved supply]
\label{def:supply}
A \emph{supply rule} is a map $\mathcal{S} : \mathbb{R}^{H} \to \mathbb{R}_{\ge 0}^{H}$
from a demand vector $d$ to a gain vector $g = \mathcal{S}(d)$ subject to the
conservation constraint
\begin{equation}
  \label{eq:conservation}
  \sum_{h=1}^{H} g_h \;=\; H .
\end{equation}
The gated layer is
\begin{equation}
  \label{eq:gated}
  \mathcal{T}_g(X) \;=\; \sum_{h=1}^{H} g_h\, W_O^{(h)} O_h .
\end{equation}
\end{definition}

Constraint \eqref{eq:conservation} is what makes the construction a \emph{supply}
problem rather than a pruning problem: total flow is redistributed, never destroyed,
and $g \equiv \mathbf{1}$ recovers the dense layer \eqref{eq:dense} exactly. All rules
below are of the form
\begin{equation}
  \label{eq:equalshare}
  g_h \;=\; \frac{H}{|S|}\,\mathbf{1}[h \in S], \qquad S \subseteq [H], \; S \neq \emptyset,
\end{equation}
that is, an equal share among a \emph{perfused set} $S$, with the rules differing only
in how $S$ is selected from $d$. We write $B \equiv |S|$ for the perfused count.

\subsection{Why the surviving circuit is the whole circuit}

The following is elementary but it is the statement that licenses mechanistic analysis
of the starved model, and it is the precise sense in which the method yields something
``small enough to verify''.

\begin{proposition}[Exact reduction]
\label{prop:reduction}
Let $g$ be of the form \eqref{eq:equalshare} with perfused set $S$, $|S| = B$. Then
$\mathcal{T}_g$ is \emph{identical}, not merely close, to a $B$-head layer:
\begin{equation}
  \mathcal{T}_g(X) \;=\; \sum_{h \in S} \widetilde{W}_O^{(h)} O_h,
  \qquad \widetilde{W}_O^{(h)} \;=\; \frac{H}{B}\, W_O^{(h)} .
\end{equation}
Moreover $\partial \mathcal{T}_g / \partial \theta_h = 0$ for every parameter $\theta_h$
belonging to a head $h \notin S$, so starved heads receive no gradient through this path
and cannot carry function.
\end{proposition}

\begin{proof}
Substituting \eqref{eq:equalshare} into \eqref{eq:gated} annihilates every term with
$h \notin S$ and scales the remainder by $H/B$; absorbing that scalar into $W_O^{(h)}$
gives the stated form, which is \eqref{eq:dense} with $H$ replaced by $B$. The gradient
claim follows since $\mathcal{T}_g$ depends on $\theta_h$ only through the product
$g_h W_O^{(h)} O_h$, and $g_h = 0$.
\end{proof}

\begin{remark}[Verification cost]
Proposition~\ref{prop:reduction} converts a claim about approximation into a claim about
identity. An analyst auditing the starved model need examine $B$ heads, not $H$, and
$\binom{B}{2}$ rather than $\binom{H}{2}$ pairwise interactions, with no residual term to
bound. At $H = 32$ and $B = 4$ this is an $8\times$ reduction in single-head hypotheses
and an $83\times$ reduction in pairwise ones. The reduction is exact only when the leak is
zero; with leak $\lambda > 0$ the discarded mass is $\lambda(H - B)$ and
Proposition~\ref{prop:reduction} holds up to a term of that order.
\end{remark}

\subsection{Relation to astrocytic normalization}

Kozachkov et al.\ construct a neuron--astrocyte network whose reading phase computes
\begin{equation}
  \label{eq:koz}
  \ell_t \;=\; \frac{1}{p}\, H\varphi(q_t) + x_t,
  \qquad
  p \;=\; \frac{1}{m} \sum_{j=1}^{N} \varphi(k_j)^{\top} \varphi(q_t),
\end{equation}
where $\varphi$ is an approximate feature map for the exponential dot-product kernel.
Because the astrocyte process activations are spatially averaged before they act
(their Eq.~8), the modulation $\widetilde{P}_{i\ell} = 1/p$ is a \emph{single scalar gain
shared by every synapse in the astrocyte's domain}, and that scalar is the reciprocal of
pooled presynaptic drive. Equation \eqref{eq:koz} is therefore a divisive gain field of
exactly the algebraic type in Definition~\ref{def:supply}: a pooled demand statistic is
inverted and applied uniformly across an index set, with the pooling enforcing a
conservation law. In their case the index set is the synapses of one astrocyte domain and
the conserved quantity is the softmax partition function; in ours the index set is the
attention heads and the conserved quantity is total flow.

This correspondence extends to the pooling parameter. Writing $\beta \in [0,1]$ for the
degree to which demand is averaged within a territory,
\begin{equation}
  d^{(\beta)}_h \;=\; (1-\beta)\, d_h \;+\; \beta \,
  \frac{1}{|\mathcal{N}(h)|} \sum_{h' \in \mathcal{N}(h)} d_{h'},
\end{equation}
the construction of \cite{kozachkov2023} is the fully pooled endpoint $\beta = 1$, in
which every element of the domain sees the same demand and hence the same gain.

\begin{remark}[Scope of the correspondence]
\label{rem:scope}
The correspondence above is \emph{algebraic and nothing more}. Kozachkov et al.\ show
that astrocytes furnish a plausible substrate for the normalization in attention; they do
not model, measure or test metabolic gating, and their astrocyte gain is a function of
presynaptic drive, not of blood supply. Citing that work as biological evidence
\emph{for} supply-based head gating would conflate substrate with mechanism. What it
legitimately supplies is the observation that a conserved divisive gain field over a
shared pool is already the operative algebra in one biologically grounded account of
attention, which motivates asking what a second such field, indexed over heads, would do.
Any physiological claim about metabolic regulation of attention requires its own
evidence, which is not offered here.
\end{remark}

\subsection{An impossibility result for fixed-quantile supply}

Let $k^\star$ denote the task's true minimal head count. The natural question is whether
$B$ discovers $k^\star$. The next proposition shows that one large family of supply rules
cannot, for reasons that have nothing to do with training.

\begin{theorem}[Fixed-quantile rules do not track $k^\star$]
\label{thm:impossible}
Let the demand be standardized, $\frac{1}{H}\sum_h d_h = 0$ and
$\frac{1}{H}\sum_h d_h^2 = 1$, and let the perfused set be
$S_\kappa = \{ h : d_h > \kappa \}$ for a fixed $\kappa$. Let $F_d$ be the empirical
distribution function of $d$. Then
\begin{equation}
  \label{eq:quantile}
  B \;=\; H\bigl(1 - F_d(\kappa)\bigr),
\end{equation}
so $B$ is a functional of $(\kappa, H)$ and the \emph{shape} of $F_d$ alone. Consequently
\begin{equation}
  \frac{\mathrm{d}B}{\mathrm{d}k^\star}
  \;=\; -H \, \frac{\partial F_d}{\partial k^\star}(\kappa),
\end{equation}
and $B$ tracks $k^\star$ only if the demand distribution deforms with $k^\star$ at
precisely the rate $\partial F_d / \partial k^\star = -1/H$ at the operating point. No
property of the rule can enforce this, because the rule has no access to $k^\star$ or to
any quantity that depends on it.
\end{theorem}

\begin{proof}
Immediate from the definition of $F_d$: the number of coordinates exceeding $\kappa$ is
$H(1 - F_d(\kappa))$. Standardization removes location and scale, so the count depends on
$d$ only through its shape. Differentiating in $k^\star$ gives the second claim.
\end{proof}

\begin{remark}[Empirical confirmation]
\label{rem:impossible}
Theorem~\ref{thm:impossible} predicts a slope near zero whenever the shape of $F_d$ is
approximately invariant across tasks, which is the generic case when $d$ is a
standardized statistic of an architecture held fixed. Sweeping the true circuit size over
$k^\star \in \{2,3,4,6,8\}$ at fixed $\kappa$ on the \texttt{multi\_relation} benchmark
gives a fitted $\mathrm{d}B/\mathrm{d}k^\star = 0.29$ against the ideal $1.00$, with mean
absolute error $1.53$ heads and $B$ saturating at $4$ for every $k^\star \ge 4$. At
$k^\star \in \{6,8\}$ the starved model fails the task outright, at losses $0.34$ and
$0.50$ against a trivial baseline of $1.007$, while passing \emph{through} the correct
perfusion level earlier in annealing with the task solved to $8\times 10^{-5}$. The rule
discards the information required to stop.
\end{remark}

\subsection{Autoregulated supply}

Theorem~\ref{thm:impossible} indicates the repair: the operating point must be set by a
quantity that depends on task difficulty. We close the loop on a metabolic deficit, in
the manner of functional hyperemia, where flow answers to a demand signal that rises when
supply is inadequate.

\begin{definition}[Autoregulated supply]
\label{def:autoreg}
Fix a target loss $L^\star$. Let $\widehat{L}_t$ be a smoothed estimate of the running
loss. The threshold is a state variable evolving as
\begin{equation}
  \label{eq:loop}
  \kappa_{t+1} \;=\; \Pi_{[\kappa_-, \kappa_+]}\Bigl(
    \kappa_t + \eta \bigl( \log L^\star - \log \widehat{L}_t \bigr) \Bigr),
\end{equation}
with $\Pi$ the projection onto the admissible interval, and $S = S_{\kappa_t}$ as before.
The control error is a log ratio, so the loop is invariant to the scale of the loss; the
target is taken as a fixed fraction of the trivial baseline, so it is invariant across
tasks.
\end{definition}

Let $L_B$ denote the smallest loss achievable with $B$ perfused heads, a nonincreasing
function of $B$. Because the head count is an integer, $\{L_B\}$ is a staircase.

\begin{theorem}[The loop selects the minimal sufficient circuit]
\label{thm:autoreg}
Suppose $L_B$ is nonincreasing in $B$ and that $\kappa \mapsto B(\kappa)$ is nonincreasing
(more ischemia perfuses no more heads). Then any fixed point of \eqref{eq:loop} in the
interior of $[\kappa_-, \kappa_+]$ satisfies $\widehat{L} = L^\star$, and the selected
perfusion is
\begin{equation}
  B^\star(L^\star) \;=\; \min \{\, B : L_B \le L^\star \,\}.
\end{equation}
In particular, if the target lies in the staircase gap at $k^\star$,
\begin{equation}
  \label{eq:gap}
  L_{k^\star} \;<\; L^\star \;<\; L_{k^\star - 1},
\end{equation}
then $B^\star(L^\star) = k^\star$, and this holds for \emph{every} $L^\star$ in that
interval.
\end{theorem}

\begin{proof}
At an interior fixed point the increment in \eqref{eq:loop} vanishes, so
$\log L^\star = \log \widehat{L}$. Monotonicity of $L_B$ in $B$ and of $B$ in $\kappa$
makes the composed map $\kappa \mapsto L_{B(\kappa)}$ nondecreasing, so the loop is a
scalar monotone feedback and its fixed point is the least $B$ meeting the target.
Condition \eqref{eq:gap} places that least $B$ at $k^\star$ by definition of the
staircase.
\end{proof}

\begin{remark}[Why this is not a tuned hyperparameter]
\label{rem:notuning}
Theorem~\ref{thm:autoreg} makes a sharp, falsifiable prediction: the selected count is
constant in $L^\star$ throughout the interval \eqref{eq:gap}, and changes only when the
target crosses a stair. The width of that interval is a property of the benchmark, not of
the method. On \texttt{multi\_relation} at $R = 4$ the measured staircase is
$L_1 = 0.755$, $L_2 = 0.503$, $L_4 = 2\times 10^{-5}$ against a trivial baseline of
$1.007$, so the admissible interval spans roughly four orders of magnitude. This
distinguishes Theorem~\ref{thm:autoreg} from the fixed-$\kappa$ rule, whose successful
setting was a single point. The corresponding control, sweeping $L^\star$ over a
$40\times$ range, is the experiment that decides whether the prediction holds.
\end{remark}

\begin{remark}[Measured behaviour]
\label{rem:autoreg}
With the loop slowed so that its time constant exceeds that of the learner, the perfused
count settles at $B = k^\star$ for every $k^\star \in \{2,3,4,6,8\}$ and every seed, with
interquartile spread $0$ or $1$ over the held phase: fitted
$\mathrm{d}B/\mathrm{d}k^\star = 1.00$ and mean absolute error $0.00$ heads, against
$0.29$ and $1.53$ for the fixed-quantile rule. At $k^\star \in \{6,8\}$, where the
fixed rule saturated and failed, the loop reaches $6$ and $8$ heads and solves the task.
Two caveats are recorded here because they bear on the interpretation. First, at
$k^\star = 3$ the count is exactly correct while the loss remains at $0.284$: the right
size is selected for a model that has not solved the task, so Theorem~\ref{thm:autoreg}
describes the selection and not the optimization. Second, role coverage does not return
to unity above $k^\star = 2$, so recovering the correct head \emph{count} is weaker than
recovering the correct head \emph{roles}. An undamped loop, with gain large enough that
the controller is faster than the learner, instead limit-cycles between $4$ and $32$
heads; the requirement is a separation of timescales, not the feedback alone.
\end{remark}

\subsection{What is and is not claimed}

Theorem~\ref{thm:autoreg} states that a conserved supply field steered by a metabolic
deficit selects the minimal head count sufficient for a target, and
Proposition~\ref{prop:reduction} states that the resulting model \emph{is} a model of
that size, exactly. Together these give a circuit that is both minimal and faithfully
analysable, which is the requirement mechanistic interpretability places on any reduction.

Three limitations are load-bearing. First, $B^\star(L^\star)$ coincides with $k^\star$ by
virtue of \eqref{eq:gap}, so the method finds the minimal \emph{sufficient} circuit rather
than discovering circuit size from metabolic principles alone; the substantive content is
that one run of a feedback loop returns what otherwise requires a sweep over $B$ and a
training run at each value. Second, Theorem~\ref{thm:impossible} is a statement about
fixed-quantile rules and does not exclude other open-loop rules that read the shape of
$F_d$. Third, and per Remark~\ref{rem:scope}, no part of this argument draws evidential
support from the astrocyte literature; the relation established in \eqref{eq:koz} is an
identity of algebraic form, and the physiological question of whether attention in
biological tissue is metabolically gated is untouched by it.

\begin{thebibliography}{9}
\bibitem{kozachkov2023}
L.~Kozachkov, K.~V.~Kastanenka, and D.~Krotov,
\emph{Building transformers from neurons and astrocytes},
Proc.\ Natl.\ Acad.\ Sci.\ U.S.A.\ \textbf{120}(34), e2219150120 (2023).
\end{thebibliography}
